Las capacidades de aislamiento y sandboxing de Piston ya se describieron en la \autoref{SS:DI_TEC_MOTOR}. Este apartado se centra en cómo se integra ese motor en el flujo de la aplicación desde que el alumno pulsa el botón de ejecución hasta que el resultado aparece en la terminal.

El alumno escribe código en el editor Monaco y pulsa ejecutar. Vue envía una petición POST autenticada al \textit{endpoint} \texttt{/compiler/execute/} de Django con el código fuente, el lenguaje y la versión del \textit{runtime}. Django construye el \textit{payload} conforme al contrato de la API interna de Piston y lo reenvía:

\PythonCode[lst:piston_call]{Llamada a Piston}{Construcción del \textit{payload} y llamada a la API de ejecución de Piston (\texttt{compiler/views.py}).}{python/p.py}{1}{8}{1}

La URL del servicio se lee de la variable de entorno \texttt{PISTON\_URL}, lo que permite apuntar a distintas instancias sin modificar el código. Piston ejecuta el programa en un contenedor efímero y devuelve un objeto JSON con los campos \texttt{compile} y \texttt{run}. Django normaliza la respuesta extrayendo \texttt{stdout}, \texttt{stderr} y el código de salida, y mapea los códigos de estado de Piston a mensajes en español orientados al alumno, de modo que valores como \texttt{TO} (timeout), \texttt{CE} (error de compilación) o \texttt{SG} (señal del sistema operativo) se traducen en explicaciones comprensibles en lugar de códigos opacos.

La respuesta normalizada llega al \textit{frontend} y se muestra en el panel de salida integrado. Pinia almacena la última salida en el estado del editor y, cuando el alumno envía el siguiente mensaje al tutor, ese valor se incluye automáticamente en la petición. Django lo inyecta como contexto adicional en el LLM, cerrando el ciclo en el que el tutor puede razonar sobre el comportamiento real del código sin que el alumno tenga que copiar y pegar ninguna traza de error.
